\scanpage{53}\mspage{- 133 -}
\hfill n\textsuperscript{o} 371

$n-2$ \uline{de} $\underline{X}$ \uline{obtenu en supprimant dans} $\underline{X}$ \uline{les lignes d'indices} $i$ \uline{et} $j$ \uline{et les colonnes d'indices} $h$ \uline{et} $k$ ; \uline{on a alors}
\[
\csname tagsleft@true\endcsname
X^{ih}X^{jk}-X^{ik}X^{jh}=(\det\underline{X}).X^{ij,hk}.\tag*{(38)}
\]
Il suffit d'appliquer la formule (35) lorsque $H=I-\{i,j\}$, $K=I-\{h,k\}$, où $I=[\![1,n]\!]$.

Nous verrons au chap.IV que les formules (34) à (38) sont encore valables sans aucune hypothèse sur la matrice carrée $\underline{X}$.

\begin{center}
-:-:-:-:-
\end{center}
% archive sheet number: 52
